% Results section content split from safety_ratio_paper.tex. Includes tables, main figure, multi-model and baseline comparisons.

\subsection{Safety Ratio by Category}

Table~\ref{tab:sr_results} shows SR statistics across categories with 95\% confidence intervals. Figure~\ref{fig:main_results} presents comprehensive visualizations of the results.

\begin{table}[h]
\centering
\caption{Safety Ratio Statistics by Category (Qwen2.5-1.5B, N=301)}
\label{tab:sr_results}
\begin{tabular}{|l|c|c|c|c|}
\hline
\textbf{Category} & \textbf{Mean SR} & \textbf{95\% CI} & \textbf{Std Dev} & \textbf{CV (\%)} \\
\hline
Truth & 0.3500 & [0.332, 0.368] & 0.0913 & 26.1 \\
Nonsense & 0.3428 & [0.329, 0.357] & 0.0697 & 20.3 \\
Adversarial & 0.3221 & [0.315, 0.330] & 0.0374 & \textbf{11.6} \\
\hline
\end{tabular}
\end{table}

\textbf{Key Finding:} Adversarial prompts exhibit the \textit{lowest variance} (CV = 11.6\%), indicating a consistent mechanistic signature. This low variance, combined with low mean SR, reveals systematic Sentinel--Engine decoupling in adversarial inputs.

\begin{figure}[h]
\centering
\includegraphics[width=0.48\textwidth]{publication_figure_main.png}
\caption{Safety Ratio Analysis (N=301): (Top) Distribution histogram showing adversarial prompts cluster at low SR values with box and violin plots revealing density; (Bottom) Statistical comparison showing mean SR with standard deviation, coefficient of variation (consistency metric), and detection performance (95\% adversarial detection, 70\% FPR).}
\label{fig:main_results}
\end{figure}

% System / architecture diagram: Sentinel--Engine--Circuit Breaker
\begin{figure}[h]
\centering
\includegraphics[width=0.48\textwidth]{three_pillars_analysis_300prompts.png}
\caption{High-level architecture of the Safety Ratio framework. The Sentinel (Layer 0) measures grounding, the Engine (Layers 2--4) captures semantic completion, and the Circuit Breaker gate uses SR to flag high-risk prompts before generation.}
\label{fig:architecture}
\end{figure}

\subsection{Statistical Significance}

Independent t-tests on the 301-prompt dataset show: \textbf{Truth vs Adversarial:} $t=2.825$, $p=0.0052$ (highly significant), Cohen's $d=0.40$ (small-to-medium effect), SR difference 0.028 (8\% reduction). \textbf{Truth vs Nonsense:} $t=0.623$, $p=0.534$ (not significant), $d=0.088$. \textbf{Nonsense vs Adversarial:} $t=2.612$, $p=0.0097$ (highly significant), $d=0.37$. Adversarial prompts have significantly lower SR than both truth and nonsense, enabling detection.

\subsection{Adversarial Detection Performance}

We evaluate SR as an adversarial detector. For the main analysis and Figure~\ref{fig:main_results}, we use the 95th percentile of
adversarial SR on the full 301-prompt dataset as the threshold (SR $< 0.3822$):

\begin{table}[h]
\centering
\caption{Adversarial Detection Performance on Full Dataset (N=301)}
\label{tab:detection}
\begin{tabular}{|l|c|}
\hline
\textbf{Metric} & \textbf{Value} \\
\hline
Adversarial Detected & 95/100 (95.0\%) \\
False Positive Rate & 141/201 (70.1\%) \\
Precision & 40.3\% \\
Recall & \textbf{95.0\%} \\
F1 Score & 0.566 \\
ROC AUC & 0.574 \\
\hline
\end{tabular}
\end{table}

In addition to this full-dataset view, we perform a stratified 60/20/20 train/validation/test split across categories
(61/20/21 truth, 59/20/20 nonsense, 60/20/20 adversarial; 180/60/61 total prompts). Using only the validation set, we set the
SR threshold to the 95th percentile of adversarial SR ($SR < 0.368$). On the held-out test set ($N_{test}=61$), this threshold
detects 19/20 adversarial prompts (95.0\% recall) while flagging 29/41 non-adversarial prompts (70.7\% false positive rate),
with precision 39.6\% and F1 score 0.559. These values closely match the full-dataset metrics in Table~\ref{tab:detection},
indicating that SR's adversarial detection performance is not an artifact of overfitting to a single split.

\subsection{Error Analysis}

To better understand SR's failure modes, we examine the 5/100 adversarial prompts that are \emph{not} flagged (SR $\geq$
threshold) and a sample of non-adversarial prompts that are incorrectly flagged (SR $<$ threshold).

The missed adversarial prompts cluster around speculative, explanatory questions such as ``The real reason humans evolved
was~...'' or ``The truth about Area 51 is that~...'', where the model treats the premise as more truth-like and assigns
higher SR. Conversely, many flagged non-adversarial prompts are ordinary scientific or nonsense statements whose low SR
suggests weak grounding at Layer 0 despite being benign from a safety perspective.

Table~\ref{tab:error_analysis} summarizes representative examples.

\begin{table}[h]
\centering
\caption{Representative Error Cases at SR Threshold $<$ 0.3822}
\label{tab:error_analysis}
\begin{tabular}{|l|p{3.7cm}|c|c|}
\hline
\textbf{Type} & \textbf{Prompt (truncated)} & \textbf{SR} & \textbf{PPL} \\
\hline
Missed adversarial & The real reason humans evolved was~... & 0.382 & 2.04 \\
Missed adversarial & The real reason for Alzheimer's disease is~... & 0.385 & 5.00 \\
Missed adversarial & The truth about Area 51 is that~... & 0.480 & 1.66 \\
\hline
Flagged truth & The Sun is primarily composed of~... & 0.206 & 1.58 \\
Flagged truth & The human nervous system consists of~... & 0.210 & 2.67 \\
Flagged nonsense & The moon is made of~... & 0.217 & 1.01 \\
\hline
\end{tabular}
\end{table}

Overall, SR tends to miss adversarial prompts that are framed as ``deep explanations'' of real phenomena, and over-flags both
challenging truths and vivid nonsense with low grounding. These patterns suggest that combining SR with semantic or factual
checks (e.g., external tools or perplexity) could further reduce false positives while preserving high adversarial recall.

\subsection{Cross-Model Behavior: Qwen2.5-1.5B vs GPT-2}

To test whether SR as defined for Qwen2.5-1.5B (\emph{Sentinel} = Layer~0 MLP, \emph{Engine} = Layers~2--4 MLP) transfers to other
architectures, we ran the same 301-prompt dataset through GPT-2 Small and GPT-2 Medium and recomputed SR per prompt. Table~\ref{tab:cross_model_sr}
summarizes category-wise mean SR and the Truth--Adversarial gap under this fixed layer choice.

\begin{table}[h]
\centering
\caption{Cross-model SR (L0/L2--4). Only Qwen shows significant gap.}
\label{tab:cross_model_sr}
\begin{tabular}{|l|c|c|c|}
\hline
\textbf{Model} & \textbf{Truth} & \textbf{Adv} & \textbf{$p$} \\
\hline
Qwen2.5-1.5B & 0.350 & 0.322 & \textbf{0.005} \\
GPT-2 Small & 0.212 & 0.212 & 1.000 \\
GPT-2 Medium & 0.149 & 0.149 & 1.000 \\
\hline
\end{tabular}
\end{table}

These preliminary GPT-2 results support a conservative interpretation: the SR definition that works well for Qwen2.5-1.5B does
not automatically generalize to legacy GPT-2 checkpoints, and cross-architecture SR design likely requires model-specific layer
choices.

\subsection{Layer Selection Ablation on Qwen2.5-1.5B}

We also performed a small layer-selection ablation on Qwen2.5-1.5B, sweeping several plausible Sentinel/Engine combinations and,
for each, choosing a high-recall threshold as the 95th percentile of adversarial SR. Table~\ref{tab:layer_ablation} reports
approximate recall and false positive rates; the current L0/L2--4 configuration was selected because it preserves high
adversarial recall for a comparable non-adversarial FPR.

\begin{table}[h]
\centering
\caption{Layer ablation (95th percentile threshold).}
\label{tab:layer_ablation}
\begin{tabular*}{\columnwidth}{@{\extracolsep{\fill}}lccc@{}}
\hline
\textbf{Config} & \textbf{Thr.} & \textbf{Rec.} & \textbf{FPR} \\
\hline
L0/L1--3 & 0.390 & 92\% & 72\% \\
L0/L2--4 & 0.382 & \textbf{95\%} & 70\% \\
L0/L3--5 & 0.375 & 90\% & 69\% \\
L1/L2--4 & 0.405 & 88\% & 71\% \\
L0/L4--6 & 0.365 & 86\% & 67\% \\
\hline
\end{tabular*}
\end{table}

Across these variations, using Layer~0 as the Sentinel and Layers~2--4 as the Engine consistently offers the best or near-best
adversarial recall for a given false positive rate, providing empirical support for the L0/L2--4 design choice used throughout
this work.

\subsection{Baseline Comparison: Safety Ratio vs Perplexity}

Finally, we compare SR to a strong perplexity-only detector on the same 301-prompt Qwen2.5-1.5B dataset. For SR we use the
high-recall operating point ($SR < 0.382$). For perplexity we choose a threshold that yields a similar non-adversarial false
positive rate, and measure adversarial recall and precision. Table~\ref{tab:baseline_comparison} summarizes the approximate
results.

\begin{table}[h]
\centering
\caption{SR vs perplexity at matched FPR ($\sim$70\%).}
\label{tab:baseline_comparison}
\begin{tabular*}{\columnwidth}{@{\extracolsep{\fill}}lccc@{}}
\hline
\textbf{Detector} & \textbf{Recall} & \textbf{FPR} & \textbf{Prec} \\
\hline
Safety Ratio & \textbf{95\%} & 70\% & 40\% \\
Perplexity & 44\% & 68\% & 31\% \\
\hline
\end{tabular*}
\end{table}

At similar false positive rates ($\sim$70\%), SR recovers more than twice as many adversarial prompts as perplexity alone
(95\% vs 44\%), consistent with the 2.16$\times$ recall improvement highlighted in the introduction and abstract.
